\section{Conclusion}
\label{sec:conclusion}

Memory-safety detectors see exactly as far as their allocator lets them. On 16 controlled heap bugs,
ASan, Valgrind Memcheck and HWASan each detected 12--15 on \code{malloc} and one or two inside an
unannotated custom allocator, whether the allocator was one of ours or one of SQLite's, and the ones
they kept were detections of the pool rather than of the object. The tools' annotation interfaces
recover most of the loss, ASan's at almost no cost and Memcheck's at about twice its already large
cost, but each interface leaves a characteristic set of bugs invisible, because poisoning describes
addressability, memory pools describe ownership and tags describe provenance at 16-byte granularity.
Rewriting the allocator away restores \code{malloc}'s detection and memory cost, at the price of the
behaviour the allocator existed for and of the bugs that depend on it. The results hold, cell for
cell, on a second machine, and every miss has a mechanism. For testing a program with a CMA, the
practical advice is to annotate the allocator, to run more than one kind of detector over it, and to
treat a clean run of an unannotated build as evidence about the code outside the pool only.

\paragraph{Team contributions.} Cheng Zefeng (Part A) designed the corpus, negative and positive
controls and the harness, ran the formal detection matrices and the HWASan study, and reviewed every
miss. Wei Yuning (Part B) implemented the allocators, their annotations and rewrites, the SQLite
allocator layers, and the EC2 matrices and allocator benchmark. Xu Ruofan (Part C) set up and pinned
the tools, investigated HWASan and LowFat feasibility, and ran the SQLite and Lua benchmarks. All
three wrote the report.
